\usepackage[english]{babel}
\usepackage[T1]{fontenc}
\usepackage{microtype}
\usepackage[english,capitalise]{cleveref}

\usepackage[acronym]{glossaries}
\makenoidxglossaries

%https://tex.stackexchange.com/questions/191376/autoref-and-braces-around-equation-number
\makeatletter
\let\oldtheequation\theequation
\renewcommand\tagform@[1]{\maketag@@@{\ignorespaces#1\unskip\@@italiccorr}}
\renewcommand\theequation{(\oldtheequation)}
\makeatother
\renewcommand{\equationautorefname}{Eq.}
\renewcommand{\figureautorefname}{Fig.}
\renewcommand{\tableautorefname}{Table}
\renewcommand{\partautorefname}{Part}
\renewcommand{\chapterautorefname}{Chapter}
\renewcommand{\sectionautorefname}{Section}
\renewcommand{\subsectionautorefname}{Section}
\renewcommand{\subsubsectionautorefname}{Section}
\renewcommand{\appendixautorefname}{Appendix}
%https://tex.stackexchange.com/questions/458671/how-to-autoref-an-appendix-in-ieee-template
\newcommand{\appref}[1]{\hyperref[#1]{Appendix~\ref*{#1}}}
\newcommand{\refcite}[1]{Ref.~\cite{#1}}

\renewcommand{\d}{\text{d}}
\renewcommand{\r}{\text{r}}
\renewcommand{\t}{\text{t}}
\renewcommand{\a}{\text{a}}
\renewcommand{\o}{\text{o}}
\renewcommand{\F}{\text{F}}
\renewcommand{\c}{\text{c}}
\renewcommand{\i}{\text{i}}
\renewcommand{\L}{\text{L}}
\renewcommand{\K}{\hat{\mathcal{K}}}
\renewcommand{\S}{\hat{\mathcal{S}}}

\newcommand{\td}{\text{td.}}
\newcommand{\U}{\hat{\mathcal{U}}}
\newcommand{\T}{\hat{\mathcal{T}}}
\newcommand{\C}{\hat{\mathcal{C}}}
\newcommand{\ret}{\text{ret}}
\newcommand{\PTR}{R}
\newcommand{\eplus}{^{\scriptscriptstyle{(+)}}}
\newcommand{\eminus}{^{\scriptscriptstyle{(-)}}}
\newcommand{\epsmall}[1]{^{\scriptscriptstyle{(#1)}}}
\newcommand{\asymeq}[1]{\overset{#1}{\simeq}}
\newcommand{\asympropto}[1]{\overset{#1}{\sim}}
%t == text commands, s == symbols
\newcommand{\tint}{\text{int}}
\newcommand{\tlin}{\text{linear}}
\newcommand{\ttherm}{\text{therm}}
\newcommand{\spage}{p.\,}
\newcommand{\spages}{pp.\,}
\newcommand{\ssec}{sec.\,}
\newcommand{\schap}{ch.\,}
\newcommand{\sapp}{app.\,}
\newcommand{\seqn}{eqn.\,}
\newcommand{\seqnl}{Equation}
\newcommand{\sappl}{Appendix}
\newcommand{\sact}{\mathcal{S}}
\newcommand{\sgf}{G}
\newcommand{\tsg}{\text{SG}}
\newcommand{\tms}{momentum-shell}
\newcommand{\tMs}{Momentum-shell}
\newcommand{\rgbeta}{\bm{\beta}}
\newcommand{\thermit}{Hermitian}
\newcommand{\trs}{\hat{\mathcal{T}}}
\newcommand{\phs}{\hat{\mathcal{C}}}
\newcommand{\chs}{\hat{\mathcal{S}}}
\newcommand{\ccg}{\hat{\mathcal{K}}}
% https://tex.stackexchange.com/questions/30619/what-is-the-best-symbol-for-vector-matrix-transpose
\newcommand{\transpose}{{\mkern-1.5mu\text{T}}}
\newcommand{\eff}{\text{eff}}
%names
\newcommand{\nameKallen}{Källén}
\newcommand{\myparagraph}[1]{{\vspace*{1em}}\,\\\textbf{#1.}\ }

% https://tex.stackexchange.com/questions/10555/hyperref-warning-token-not-allowed-in-a-pdf-string
\pdfstringdefDisableCommands{%
    \def\\{}%
}
% glossary references not coloured (https://tex.stackexchange.com/questions/38544/glossary-links-color)
\makeatletter
\newcommand*{\glsplainhyperlink}[2]{%
    \begingroup%
      \hypersetup{hidelinks}%
      \hyperlink{#1}{#2}%
    \endgroup%
}
\let\@glslink\glsplainhyperlink
\makeatother
% https://tex.stackexchange.com/questions/50747/options-for-appearance-of-links-in-hyperref
%\definecolor{mDarkBlue}{rgb}{0,0,0.4}
\hypersetup{
    colorlinks,
    breaklinks,
    citecolor=blue,
    linkcolor=blue,
    urlcolor=blue,
    pdfauthor={\Author},
    pdftitle={\Title}
}